% Options for packages loaded elsewhere
% Options for packages loaded elsewhere
\PassOptionsToPackage{unicode}{hyperref}
\PassOptionsToPackage{hyphens}{url}
\PassOptionsToPackage{dvipsnames,svgnames,x11names}{xcolor}
%
\documentclass[
  letterpaper,
  DIV=11,
  numbers=noendperiod]{scrartcl}
\usepackage{xcolor}
\usepackage{amsmath,amssymb}
\setcounter{secnumdepth}{5}
\usepackage{iftex}
\ifPDFTeX
  \usepackage[T1]{fontenc}
  \usepackage[utf8]{inputenc}
  \usepackage{textcomp} % provide euro and other symbols
\else % if luatex or xetex
  \usepackage{unicode-math} % this also loads fontspec
  \defaultfontfeatures{Scale=MatchLowercase}
  \defaultfontfeatures[\rmfamily]{Ligatures=TeX,Scale=1}
\fi
\usepackage{lmodern}
\ifPDFTeX\else
  % xetex/luatex font selection
\fi
% Use upquote if available, for straight quotes in verbatim environments
\IfFileExists{upquote.sty}{\usepackage{upquote}}{}
\IfFileExists{microtype.sty}{% use microtype if available
  \usepackage[]{microtype}
  \UseMicrotypeSet[protrusion]{basicmath} % disable protrusion for tt fonts
}{}
\makeatletter
\@ifundefined{KOMAClassName}{% if non-KOMA class
  \IfFileExists{parskip.sty}{%
    \usepackage{parskip}
  }{% else
    \setlength{\parindent}{0pt}
    \setlength{\parskip}{6pt plus 2pt minus 1pt}}
}{% if KOMA class
  \KOMAoptions{parskip=half}}
\makeatother
% Make \paragraph and \subparagraph free-standing
\makeatletter
\ifx\paragraph\undefined\else
  \let\oldparagraph\paragraph
  \renewcommand{\paragraph}{
    \@ifstar
      \xxxParagraphStar
      \xxxParagraphNoStar
  }
  \newcommand{\xxxParagraphStar}[1]{\oldparagraph*{#1}\mbox{}}
  \newcommand{\xxxParagraphNoStar}[1]{\oldparagraph{#1}\mbox{}}
\fi
\ifx\subparagraph\undefined\else
  \let\oldsubparagraph\subparagraph
  \renewcommand{\subparagraph}{
    \@ifstar
      \xxxSubParagraphStar
      \xxxSubParagraphNoStar
  }
  \newcommand{\xxxSubParagraphStar}[1]{\oldsubparagraph*{#1}\mbox{}}
  \newcommand{\xxxSubParagraphNoStar}[1]{\oldsubparagraph{#1}\mbox{}}
\fi
\makeatother


\usepackage{longtable,booktabs,array}
\usepackage{calc} % for calculating minipage widths
% Correct order of tables after \paragraph or \subparagraph
\usepackage{etoolbox}
\makeatletter
\patchcmd\longtable{\par}{\if@noskipsec\mbox{}\fi\par}{}{}
\makeatother
% Allow footnotes in longtable head/foot
\IfFileExists{footnotehyper.sty}{\usepackage{footnotehyper}}{\usepackage{footnote}}
\makesavenoteenv{longtable}
\usepackage{graphicx}
\makeatletter
\newsavebox\pandoc@box
\newcommand*\pandocbounded[1]{% scales image to fit in text height/width
  \sbox\pandoc@box{#1}%
  \Gscale@div\@tempa{\textheight}{\dimexpr\ht\pandoc@box+\dp\pandoc@box\relax}%
  \Gscale@div\@tempb{\linewidth}{\wd\pandoc@box}%
  \ifdim\@tempb\p@<\@tempa\p@\let\@tempa\@tempb\fi% select the smaller of both
  \ifdim\@tempa\p@<\p@\scalebox{\@tempa}{\usebox\pandoc@box}%
  \else\usebox{\pandoc@box}%
  \fi%
}
% Set default figure placement to htbp
\def\fps@figure{htbp}
\makeatother


% definitions for citeproc citations
\NewDocumentCommand\citeproctext{}{}
\NewDocumentCommand\citeproc{mm}{%
  \begingroup\def\citeproctext{#2}\cite{#1}\endgroup}
\makeatletter
 % allow citations to break across lines
 \let\@cite@ofmt\@firstofone
 % avoid brackets around text for \cite:
 \def\@biblabel#1{}
 \def\@cite#1#2{{#1\if@tempswa , #2\fi}}
\makeatother
\newlength{\cslhangindent}
\setlength{\cslhangindent}{1.5em}
\newlength{\csllabelwidth}
\setlength{\csllabelwidth}{3em}
\newenvironment{CSLReferences}[2] % #1 hanging-indent, #2 entry-spacing
 {\begin{list}{}{%
  \setlength{\itemindent}{0pt}
  \setlength{\leftmargin}{0pt}
  \setlength{\parsep}{0pt}
  % turn on hanging indent if param 1 is 1
  \ifodd #1
   \setlength{\leftmargin}{\cslhangindent}
   \setlength{\itemindent}{-1\cslhangindent}
  \fi
  % set entry spacing
  \setlength{\itemsep}{#2\baselineskip}}}
 {\end{list}}
\usepackage{calc}
\newcommand{\CSLBlock}[1]{\hfill\break\parbox[t]{\linewidth}{\strut\ignorespaces#1\strut}}
\newcommand{\CSLLeftMargin}[1]{\parbox[t]{\csllabelwidth}{\strut#1\strut}}
\newcommand{\CSLRightInline}[1]{\parbox[t]{\linewidth - \csllabelwidth}{\strut#1\strut}}
\newcommand{\CSLIndent}[1]{\hspace{\cslhangindent}#1}



\setlength{\emergencystretch}{3em} % prevent overfull lines

\providecommand{\tightlist}{%
  \setlength{\itemsep}{0pt}\setlength{\parskip}{0pt}}



 


\KOMAoption{captions}{tableheading}
\makeatletter
\@ifpackageloaded{caption}{}{\usepackage{caption}}
\AtBeginDocument{%
\ifdefined\contentsname
  \renewcommand*\contentsname{Table of contents}
\else
  \newcommand\contentsname{Table of contents}
\fi
\ifdefined\listfigurename
  \renewcommand*\listfigurename{List of Figures}
\else
  \newcommand\listfigurename{List of Figures}
\fi
\ifdefined\listtablename
  \renewcommand*\listtablename{List of Tables}
\else
  \newcommand\listtablename{List of Tables}
\fi
\ifdefined\figurename
  \renewcommand*\figurename{Figure}
\else
  \newcommand\figurename{Figure}
\fi
\ifdefined\tablename
  \renewcommand*\tablename{Table}
\else
  \newcommand\tablename{Table}
\fi
}
\@ifpackageloaded{float}{}{\usepackage{float}}
\floatstyle{ruled}
\@ifundefined{c@chapter}{\newfloat{codelisting}{h}{lop}}{\newfloat{codelisting}{h}{lop}[chapter]}
\floatname{codelisting}{Listing}
\newcommand*\listoflistings{\listof{codelisting}{List of Listings}}
\usepackage{amsthm}
\theoremstyle{plain}
\newtheorem{lemma}{Lemma}[section]
\theoremstyle{definition}
\newtheorem{definition}{Definition}[section]
\theoremstyle{plain}
\newtheorem{theorem}{Theorem}[section]
\theoremstyle{plain}
\newtheorem{corollary}{Corollary}[section]
\theoremstyle{remark}
\AtBeginDocument{\renewcommand*{\proofname}{Proof}}
\newtheorem*{remark}{Remark}
\newtheorem*{solution}{Solution}
\newtheorem{refremark}{Remark}[section]
\newtheorem{refsolution}{Solution}[section]
\makeatother
\makeatletter
\makeatother
\makeatletter
\@ifpackageloaded{caption}{}{\usepackage{caption}}
\@ifpackageloaded{subcaption}{}{\usepackage{subcaption}}
\makeatother
\makeatletter
\@ifpackageloaded{academicons}{}{\usepackage{academicons}}
\makeatother
\usepackage{bookmark}
\IfFileExists{xurl.sty}{\usepackage{xurl}}{} % add URL line breaks if available
\urlstyle{same}
\hypersetup{
  pdftitle={Category Theory},
  pdfauthor={Rong-Kang Zhang},
  pdfkeywords={Natural Transformation},
  colorlinks=true,
  linkcolor={blue},
  filecolor={Maroon},
  citecolor={Blue},
  urlcolor={Blue},
  pdfcreator={LaTeX via pandoc}}


\title{Category Theory}
\author{Rong-Kang Zhang}
\date{2025-11-09}
\begin{document}
\maketitle

\renewcommand*\contentsname{Contents}
{
\hypersetup{linkcolor=}
\setcounter{tocdepth}{6}
\tableofcontents
}

\section{Category theory}\label{category-theory}

See (Lane 1998)

\subsection{Introduction}\label{introduction}

\subsubsection{Natural transformation}\label{natural-transformation}

See \textbf{?@eq-gauss} and \textbf{?@thm-fubini}.

\phantomsection\label{fubini}
Let \(f \in L^1(X\times Y)\). Then
\[ \int_X\!\int_Y f(x,y)\,dy\,dx = \int_Y\!\int_X f(x,y)\,dx\,dy. \]

\begin{definition}[]\protect\hypertarget{def-group}{}\label{def-group}

A \textbf{group} is a set \(G\) with a binary operation \(\cdot\)
satisfying: 1. Associativity: \((a\cdot b)\cdot c=a\cdot(b\cdot c)\); 2.
Identity: \(\exists e\in G\) such that \(e\cdot a=a\cdot e=a\); 3.
Inverses: \(\forall a\in G\ \exists a^{-1}\in G\) such that
\(a\cdot a^{-1}=a^{-1}\cdot a=e\).

\end{definition}

\phantomsection\label{ex-integers}
The set \(\mathbb Z\) with addition is a group.

\begin{lemma}[]\protect\hypertarget{lem-product}{}\label{lem-product}

The product of two integers is even if at least one factor is even.

\end{lemma}

\begin{theorem}[]\protect\hypertarget{thm-primes}{}\label{thm-primes}

There are infinitely many prime numbers.

\end{theorem}

\begin{proof}
Assume finitely many primes \(p_1,\dots,p_n\).\\
Consider \(P=p_1\cdots p_n+1\).\\
\(P>1\) has a prime divisor \(p\), but \(p\) divides neither \(P-1\) nor
\(1\), contradiction.
\end{proof}

\begin{corollary}[]\protect\hypertarget{cor-next-prime}{}\label{cor-next-prime}

For every integer \(n\ge 2\) there exists a prime \(p>n\).

\end{corollary}

\begin{refremark}
Euclid's proof is constructive: it gives an upper bound on the next
prime.

\label{rem-euclid}

\end{refremark}

Definition (Index set \(\left(A_i\right)_{i \in I}\) )
\(A: I \rightarrow X, i \mapsto A_i\) Definition (Product over the index
set \(\left(A_i\right)_{i \in I}\) ) \[
\prod_{i \in I} A_i:=\left\{a: I \rightarrow \bigcup_{i \in I} A_i \mid \forall i \in I, a(i) \in A_i\right\}
\]

Definition (Semigroup \((S, *)\) )

\[
*: S \times S \rightarrow S
\]

\begin{itemize}
\tightlist
\item
  Axiom 1 (Closure)
\end{itemize}

\[
\forall a, b \in S, a * b \in S
\]

\begin{itemize}
\tightlist
\item
  Axiom 2 (Associativity)
\end{itemize}

\[
\forall a, b, c \in S,(a * b) * c=a *(b * c)
\]

Example (Semigroup)

\[
(\mathbb{N},+), \mathbb{N}=\{1,2,3, \ldots\}, 0 \notin \mathbb{N}
\]

Definition (Commutative semigroup)

\[
\forall a, b \in S, a * b=b * a
\]

Definition (Non-commutative semigroup)

\[
\left\{\begin{array}{l}
\forall a, b, c \in S,(a * b) * c=a *(b * c) \\
\exists x, y \in S, x * y \neq y * x
\end{array}\right.
\] 查看是否更新

\aiicon{academia} \aiicon{academia-square} \aiicon{acclaim}
\aiicon{acclaim-square} \aiicon{acm} \aiicon{acm-square} \aiicon{acmdl}
\aiicon{acmdl-square} ★ \aiicon{ads} ★ \aiicon{ads-square}
\aiicon{africarxiv} \aiicon{africarxiv-square} \aiicon{archive}
\aiicon{archive-square} ★ \aiicon{arxiv} ★ \aiicon{arxiv-square}
\aiicon{biorxiv} \aiicon{biorxiv-square} \aiicon{ceur}
\aiicon{ceur-square} \aiicon{cienciavitae} \aiicon{cienciavitae-square}
\aiicon{clarivate} \aiicon{clarivate-square} \aiicon{closedaccess}
\aiicon{closedaccess-square} \aiicon{conversation}
\aiicon{conversation-square} \aiicon{coursera} \aiicon{coursera-square}
\aiicon{crossref} \aiicon{crossref-square} \aiicon{cv}
\aiicon{cv-square} \aiicon{datacite} \aiicon{datacite-square}
\aiicon{dataverse} \aiicon{dataverse-square} ★ \aiicon{dblp} ★
\aiicon{dblp-square} \aiicon{depsy} \aiicon{depsy-square} ★ \aiicon{doi}
★ \aiicon{doi-square} \aiicon{dryad} \aiicon{dryad-square}
\aiicon{elsevier} \aiicon{elsevier-square} \aiicon{figshare}
\aiicon{figshare-square} ★ \aiicon{googlescholar} ★
\aiicon{googlescholar-square} \aiicon{hal} \aiicon{hal-square}
\aiicon{hypothesis} \aiicon{hypothesis-square} \aiicon{ideasrepec}
\aiicon{ideasrepec-square} \aiicon{ieee} \aiicon{ieee-square}
\aiicon{impactstory} \aiicon{impactstory-square} \aiicon{inaturalist}
\aiicon{inaturalist-square} \aiicon{inpn} \aiicon{inpn-square} ★
\aiicon{inspire} ★ \aiicon{inspire-square} \aiicon{isidore}
\aiicon{isidore-square} ★ \aiicon{jstor} ★ \aiicon{jstor-square}
\aiicon{lattes} \aiicon{lattes-square} ★ \aiicon{mathoverflow} ★
\aiicon{mathoverflow-square} ★ \aiicon{mendeley} ★
\aiicon{mendeley-square} \aiicon{moodle} \aiicon{moodle-square}
\aiicon{mtmt} \aiicon{mtmt-square} \aiicon{nakala}
\aiicon{nakala-square} \aiicon{obp} \aiicon{obp-square}
\aiicon{openaccess} \aiicon{openaccess-square} \aiicon{opendata}
\aiicon{opendata-square} \aiicon{openmaterials}
\aiicon{openmaterials-square} \aiicon{openedition}
\aiicon{openedition-square} ★ \aiicon{orcid} ★ \aiicon{orcid-square} ★
\aiicon{osf} ★ \aiicon{osf-square} ★ \aiicon{overleaf} ★
\aiicon{overleaf-square} \aiicon{philpapers} \aiicon{philpapers-square}
\aiicon{piazza} \aiicon{piazza-square} \aiicon{preregistered}
\aiicon{preregistered-square} \aiicon{protocols}
\aiicon{protocols-square} \aiicon{psyarxiv} \aiicon{psyarxiv-square} ★
\aiicon{publons} ★ \aiicon{publons-square} ★ \aiicon{pubmed} ★
\aiicon{pubmed-square} ★ \aiicon{pubpeer} ★ \aiicon{pubpeer-square} ★
\aiicon{researcherid} ★ \aiicon{researcherid-square} ★
\aiicon{researchgate} ★ \aiicon{researchgate-square} \aiicon{ror}
\aiicon{ror-square} \aiicon{scihub} \aiicon{scihub-square} ★
\aiicon{scirate} ★ \aiicon{scirate-square} ★ \aiicon{scopus} ★
\aiicon{scopus-square} ★ \aiicon{semanticscholar} ★
\aiicon{semanticscholar-square} ★ \aiicon{springer} ★
\aiicon{springer-square} ★ \aiicon{ssrn} ★ \aiicon{ssrn-square} ★
\aiicon{stackoverflow} ★ \aiicon{stackoverflow-square} ★ \aiicon{zenodo}
★ \aiicon{zenodo-square} ★ \aiicon{zotero} ★ \aiicon{zotero-square}

\phantomsection\label{refs}
\begin{CSLReferences}{1}{0}
\bibitem[\citeproctext]{ref-maclane1998}
Lane, S. M. (1998), \emph{Categories for the working mathematician},
Graduate texts in mathematics, New York, NY: Springer-Verlag.
\url{https://doi.org/10.1007/978-1-4757-4721-8}.

\end{CSLReferences}




\end{document}
